\documentclass[11pt]{article}
\usepackage[a4paper,margin=26mm]{geometry}
\usepackage{amsmath,amssymb,mathtools}
\usepackage{unicode-math}
\setmainfont{Latin Modern Roman}
\setmathfont{Latin Modern Math}
\usepackage{xurl,hyperref}
\hypersetup{hidelinks,pdftitle={Problem statement}}
\newcommand{\sref}[2]{\hyperlink{source:#1:#2}{[S#2]}}
\newcommand{\cataloguescope}[1]{#1}
\begin{document}
\section{Hyperbolic geometry — Question 10.3}\label{statement}

\subsection{Definitions and mathematical statement}

\paragraph{Shared notation.}$\mathbb N=\{1,2,\ldots\}$, and $\mathbb Z,\mathbb Q,\mathbb R,\mathbb C$ denote the integers, rationals, reals and complex numbers. Any other conventions are specified locally.


A codimension-one foliation $\mathcal F$ decomposes a $3$-manifold $M$ into surface leaves, locally as the horizontal slices of a product chart. In the terminology of the source, it is taut if a closed curve transverse to $\mathcal F$ meets every leaf. A minimal foliation means that every leaf is dense; a minimal surface means a surface with vanishing mean curvature for the specified Riemannian metric. These two uses of minimal are distinct. For a finite-depth foliation, depth-zero leaves are compact; a leaf has depth $d>0$ when its frontier is a union of leaves of smaller depth and $d$ is the smallest such finite bound. The depth of the foliation is the greatest leaf depth. For a knot, depth is the minimum depth of a taut foliation of its exterior with the standard boundary behavior used for knot-depth foliations; fibered knots have depth zero. For a lifted leaf $\lambda\subset\mathbb H^3$, its limit set $\lambda_\infty$ is the intersection of its closure in $\mathbb H^3\cup S^2_\infty$ with $S^2_\infty$. Hausdorff dimension uses the round sphere metric, or an equivalent visual metric in the usual hyperbolic normalization.
\paragraph{Statement recorded in the supplied source.}
For a finite-depth foliation $\mathcal F$ of a hyperbolic $3$-manifold, determine whether, and how, $\dim_H(\lambda_\infty)$ for a lifted leaf $\lambda$ is related to the depth of its underlying leaf or to the depth of $\mathcal F$. Do not restrict to compact leaves or depth one. \sref{statement}{2}
\subsection{Short English statement}
Relate the dimension of a lifted leaf limit set to leaf depth and total foliation depth.
\subsection{Sources}
\begin{description}
\item[\hypertarget{source:statement:1}{S1}] ULAM.AI and the UnsolvedMath Contributors, \emph{UnsolvedMath: A Curated Collection of Open Mathematics Problems}, catalogue entry. CC BY 4.0. \url{https://huggingface.co/datasets/ulamai/UnsolvedMath}.
\item[\hypertarget{source:statement:2}{S2}] Danny Calegari, \emph{Problems in foliations and laminations of 3-manifolds}, version 0.78 (2002), Question 10.3; Definitions 1.1--1.10 and local definitions and remarks. \url{https://arxiv.org/abs/math/0209081v1}.
\end{description}
\label{end:statement}
\end{document}
